\documentclass[10pt]{article}
\usepackage[margin=21mm]{geometry}
\usepackage{amsmath,amssymb,amsthm,mathtools,hyperref}
\newtheorem{theorem}{Theorem}
\newtheorem{lemma}[theorem]{Lemma}
\newtheorem{proposition}[theorem]{Proposition}
\newcommand{\eps}{\varepsilon}
\newcommand{\norm}[1]{\lVert#1\rVert_{\rm row}}
\allowdisplaybreaks
\title{An explicit bounded-$R$ bridge for the complete C1 certificate}
\author{}\date{23 October 2026}
\begin{document}\maketitle

The source formulas, full parent, and root identities below are those of
\texttt{latex/certificates/c1\_short.tex}, (P1), (BE3)--(BE9), (SP3),
(SP6), and the twelve- and six-root exports in
\texttt{extracted-lifts/local-order669/}, at commit
\texttt{fb58c232e374dadf3cbb5f23e7e7413d429de884}.
The new result is a quantitative extension of their source admission
from $z\le1$ to every specified finite $R$. No coefficient map or
previously evaluated pencil is regenerated.

\begin{theorem}[Explicit bounded-$R$ admission]\label{thm:bridge}
For a real $R\ge1$, put
\[
 B=1+R,\qquad \eps_R=2^{-40}B^{-8},\qquad
 \delta_R=2^{-20}B^{-2}.
\]
Take exactly the native source formulas (P1), with
\[
 0<x,\sigma\le\eps_R,\quad 0<z\le R,\quad
 0<b,h<1,\quad r=bz<1,\quad E,M,J>0,\quad D,\beta\ge0.
\]
The complete parent $F_w$ has a separated source/testing sum of squares
on the independent rectangle
\[
             \sqrt{|u|},\ |w|\le\delta_R
\]
for every real state $(W,X,Y,Z)$ and unrestricted $H=h/\sigma$.
The same algebraic root identities and feature supports apply, including
both state normals, both circuit-sign sectors, $D=0$, and $\beta=0$.

One fixed family for each $R$ has at most $220$ functions on $\lambda>0$
and $157$ on $\lambda\le0$. Their source coefficients are given by the
matrix splits below and the already evaluated retained-root formulas.
Every numerical bound here is a real inequality, so the assertion
holds over real closed source and testing fields as well.
\end{theorem}

The theorem does not assert that the functions are independent of $R$.
It covers positive opening variables; a literal zero source opening
requires its own chart. These are the same restrictions as the native
bounded-positive-$z$ C1 chart.

\section{Source estimates that retain the exact identities}

Write $Q=1-b$, $\theta=1-h$, $e=xh/\sigma$, $E=1-e$, $K=1+x$,
$R_s=1+xr$, $M=R_s-xz\theta$, and
$J=1+xh-\sigma(1-r)$. The positive denominator is
\[
 P=\sigma+K+\sigma z(1+bx+2xh)
                  +\sigma bx^2hz^2+x^2h\theta z.
\]
The three opening-cleared coefficients are
\[
 p={z(1+xh)E+(1-r)(1+xhz)\over2P},\quad
 l={(1+\sigma r)\{eR_s(1+\sigma z)+ME\}\over2P},\quad
 \kappa={K(1+\sigma r)(1+\sigma z)\over2MP}.
\]
As in (P1), $m=\kappa xze$, $a=p/b$, $\ell=l/Q$, and
$C(1-r)=Kz(1+\sigma r)E/(2P)$; $D$ and $\beta$ are the native
nonnegative rational coefficients, not independently chosen variables.
Put
\[
 H_0=1+zQ,\qquad B_s=P+xz(J-\sigma),
\]
\[
 \mathcal H=R_sH_0+z\{(1-r)xh(2+xr)
       -e[2+x(1+h+r\theta)]+Ke^2\}.
\]
Throughout the proof $\eps$ can be any number with $0<\eps\le\eps_R$
that bounds both $x$ and $\sigma$. In particular $B^j\eps$ is small
for every exponent $j\le8$ used below.

The exact source identities remain
\begin{align}
 2\theta P\beta&=hB_s-\sigma H_0,\tag{BE3}\\
 2P(JD+x\theta\beta)&=\sigma^2\mathcal H,\tag{BE4}\\
 eB_s&=xH_0+\sigma\mathcal H-2PJD/\sigma
                      \le xH_0+\sigma\mathcal H.\label{eq:e}
\end{align}
No rescaling of the source or testing arguments has changed these equations.
The following conservative bounds suffice:
\begin{align}
 &1\le P\le1+4B\eps,\quad |J-1|\le\eps,\quad
 |M-1|\le B\eps,\quad1\le B_s\le1+6B\eps,
 \quad1\le H_0\le B,\label{eq:b1}\\
 &e\le5B\eps,\quad |\mathcal H-H_0|\le32B^2\eps,
 \quad0\le D\le2B\sigma^2,\quad0\le\beta\theta\le h,
                                                        \label{eq:b2}\\
 &|p-H_0/2|\le8B^2\eps,\quad |l-1/2|\le4B\eps,
 \quad |\kappa-1/2|\le6B\eps,\label{eq:b3}\\
 &\tfrac13\le p\le B,\quad\tfrac13\le l,\kappa\le1,
 \quad m\le5B^2\eps^2,\quad C(1-r)\le z\le B.\label{eq:b4}
\end{align}
To check these bounds, use $r<1$ before estimating the positive
summands of $P$. This gives \eqref{eq:b1}. Initially $E>0$ gives
$e<1$, and dropping the negative term in $\mathcal H$ gives
$\mathcal H\le4B$. Equation \eqref{eq:e} then gives
$e\le Bx+4B\sigma\le5B\eps$. Direct expansion yields
\[
 |\mathcal H-H_0|\le4B\eps+3Be+2Be^2\le32B^2\eps.
\]
Now $\mathcal H\le2B$, and (BE4) gives the bound on $D$; (BE3)
gives $\beta\theta\le h$. For the bounds on $p,l$, the useful exact
numerator differences are
\[
 2Pp-H_0=xzh(2-r-e)-ze,
\]
\[
 {2Pl\over1+\sigma r}=M+e\{xz\theta+\sigma zR_s\}.
\]
They prove \eqref{eq:b3}--\eqref{eq:b4} by \eqref{eq:b1}; the
formula for $\kappa$ is handled in the same way.

Retain the source definitions (SP3):
\[
 \nu=(2D-2\beta\theta+h)/\sigma,\quad
 \lambda=(8Dh+4\beta\theta^2+2h^2)/\sigma-4bp\sigma z,
 \quad\lambda_0=\lambda-2H,
\]
\[
 K_A=4\kappa H-\lambda,\quad B_x=8\kappa e,\quad
 R_0=-4\{2(C+D)(r-1)+2p+r-1\}.
\]
Then
\begin{align}
 &|\nu|\le2B,\quad |\lambda_0|\le3B,
 \quad |\lambda_0+2\theta H_0|\le64B^2\eps,\label{eq:b5}\\
 &|K_A-(2\theta H_0-2h(1-r))|\le1024B^4\eps,\label{eq:b6}\\
 &|R_0|\le32B(x+\sigma),\quad |B_x|\le40B\eps,
 \quad |R_x|\le16,\quad |A_y|\le32B,
 \quad |\chi|\le64B,\quad 0\le k_1\le16.\label{eq:b7}
\end{align}
Here $R_x,A_y,\chi,k_1$ have precisely their printed (SP3) definitions.
For \eqref{eq:b5}, use the exact rearrangements
\[
 \nu=2D/\sigma+[H_0-ez(J-\sigma)]/P,
\]
\[
 \lambda_0=2e\theta z(J-\sigma)/P-2\theta H_0/P
                                      +8Dh/\sigma-4pr\sigma.
\]
The only apparently unbounded factor in \eqref{eq:b6} is removed by
\[
 K(1+\sigma r)(1+\sigma z)-MP=xF_x+\sigma F_\sigma,
 \qquad F_\sigma=r-1+\sigma rz,
\]
\[
 H(2\kappa-1)={eF_x+hF_\sigma\over MP}.
\]
The polynomial $F_x$, after $r=bz$, has coefficient $\ell^1$ norm
$37$ and $z$-degree $3$. Thus $|F_x|\le37B^3$.
Also $(MP)^{-1}\le2$ and $|(MP)^{-1}-1|\le12B\eps$, giving
\[
 |H(2\kappa-1)+h(1-r)|\le384B^4\eps.
\]
Together with $K_A=2H(2\kappa-1)-\lambda_0$, this proves
\eqref{eq:b6}. The pole-free identity
\[
\begin{split}
 -PR_0/4={}&zE\{1+xh-K(1+\sigma r)\}\\
 &+(1-r)(1+xhz-P)-2DP(1-r)
\end{split}
\]
gives the first bound in \eqref{eq:b7}; the others follow by
substitution into (SP3), using $|h-Q|\le1$ and
$\beta\theta^2\le h\theta\le1/4$.

\paragraph{The downstream source identities.}
In particular $1/2\le\mathcal H\le2B$. Therefore the unchanged
sphere definitions
\[
 d_0^2={2PJD\over\sigma^2\mathcal H},\qquad
 b_0^2={2Px\theta\beta\over\sigma^2\mathcal H}
\]
give $d_0,b_0\ge0$, $d_0^2+b_0^2=1$ without dividing by either
coordinate, and
\[
 {\sigma^2\over8}\le JD+x\theta\beta\le B\sigma^2.
\]
The exact (BE6)--(BE8) formulas for $D,\beta,H,m_0,m_1$ are unchanged.
For example
\[
 H={H_0\over B_s}+{\sigma\over x}{\mathcal H\over B_s}b_0^2,
 \quad {1\over6}\le{m_0\over x^2z}\le B,
 \quad {1\over12}\le{m_1\over x\sigma z}\le2B.
\]
The three (BE9) errors are bounded, respectively, by
$8B\eps,8B^2\eps,4B\eps$. Finally $2P(p-l)/z$ is a polynomial
in $(x,\sigma,z,b,h,e)$ with coefficient norm $18$ and $z$-degree
$2$. Hence $|(p-l)/z|\le9B^2$. This proves all the boundedness and
asymptotic assertions used downstream, including the radial boundary
estimate, from the same exact source identities.

\section{A uniform bound for the positive-angular Gram}

Suppose $\lambda>0$. Use the existing twelve-root column
\[
 (N_x,N_y,R,\Gamma_*,A_\theta,T_\theta,d_1,d_2,d_3,d_4,U_*,K_*)^T
\]
and its exact matrix $Q(s,w)$, where $s=\sqrt u$. Put
\[
 d_G=4D/\sigma^2,\quad\rho=\lambda x/\sigma,\quad S=d_G+\rho,
 \quad L_0=\lambda\sigma,\quad q_0=\lambda\sqrt x,
 \quad c=\max(0,1-K_A),\quad a_0=\sqrt x/\sigma.
\]
The source identities and \eqref{eq:b1}--\eqref{eq:b7} imply
\begin{gather}
 |S-2H_0|\le128B^2\eps,\quad1\le S\le3B,
 \quad0\le d_G,\rho\le3B,\quad0<L_0\le3,\label{eq:sp}\\
 {\sqrt{d_G}\over S},{\sqrt\rho\over S},{d_G\over S},
 {\rho\over S}\le1,\qquad q_0/S\le2,\quad\sqrt{q_0}\le2B.
                                                        \label{eq:ratios}
\end{gather}
Indeed the exact identity to use is
\[
 S=2\mathcal H/P+d_G(1-J+2hx)+2eH_0/P
                     -2e^2z(J-\sigma)/P-4prx.
\]
Initially $d_G\le8B$; bounding the five terms gives
\eqref{eq:sp}. Then $q_0^2=L_0\rho$ gives \eqref{eq:ratios}.
Further,
\[
 (x+\sigma)^2H=xe+2\sigma e+\sigma h\le2\eps,
 \qquad |R_0\sqrt\lambda|\le128B\sqrt\eps.
\]
The second conclusion follows from $\lambda\le2H+3B$ and
$(x+\sigma)^2\lambda\le5\eps$.

The compensation terms satisfy
\begin{equation}
 0\le c\le4,\quad ca_0\le32B,\quad c/L_0\le32,
                         \quad c/\lambda\le32\sigma.\label{eq:c}
\end{equation}
For $c>0$, \eqref{eq:b6} forces $h>1/8$, because the reference
$2\theta H_0-2h(1-r)$ is at least $3/2$ when $h\le1/8$.
Now $L_0\ge h>1/8$ and
\[
 a_0^2(hB_s-\sigma H_0)\le\mathcal H\le2B,
 \qquad hB_s-\sigma H_0\ge1/16.
\]
Thus $a_0\le8B$. If $c=0$, every compensation term is zero;
no unused denominator is bounded away from zero.

Here is a finite positive-coefficient check of every matrix entry.
Let $\tau=\sqrt\eps$, and regard the twelve root letters as independent
nonnegative formal variables, abbreviating
$G=\Gamma_*$, $A=A_\theta$, $T=T_\theta$, $U=U_*$, $K'=K_*$.
The two full normal conductors $\mathfrak U=u\Gamma$ and
$\mathfrak V=(R_0-8pu^2)qw\Gamma$ have coefficient majorants
\[
 f_1=uG+\tau d_1+ud_2+(u+\tau)d_3+\tau T,
\]
\[
\begin{split}
 f_2={}&64B^2\tau^2wG+64B^2\tau^4R
 +64B^2\tau^2(w+\tau)(d_2+d_3)\\
 &+(128B^2\tau^2+128B^2\tau w)T+8B^2uwf_1.
\end{split}
\]
These follow by substituting \eqref{eq:ratios} in the two exact
conductor formulas in \texttt{twelve\_root.tex}; in particular no
normal coupling is deleted. Let $B_{\rm src}$ denote the original
source reference with $N_x$ entry $2$, radial form
$J_C^2+J_\ell^2$, $G$ entry $1$, angular entries $(K_A+c)/2$,
four circuit entries $1/2$, and final entries $1,1/2$.
Then $B_{\rm src}\succeq I_{12}/3$.
The correction $Q-B_{\rm src}$ is coefficientwise dominated by
\begin{align*}
 \mathcal P={}&(\tfrac{131}{2}B\tau^2+4u^2)R^2
 +(64B^2\tau^2+16Bu^2+8w^2)N_yR\\
 &+8BN_xf_1+(N_y+R)f_2+8BuwN_xN_y+16BuwN_xR\\
 &+40B^2\tau^2wN_xA+32BuN_yA+64BuAR+2BsUR\\
 &+32BuU^2+64BsUK'+32B\tau(K')^2\\
 &+[(64B+32)u+32\tau]G^2
 +(64\tau^2+4u+4\tau)\sum_i d_i^2.
\end{align*}
For example, $|E_{RR}^0|\le64B\eps$,
the coefficients of $E_{yR}$ on $(1,u^2,w^2)$ are at most
$(64B^2\eps,16B,8)$, and the three compensation expressions
are bounded using \eqref{eq:c}. These give exactly the displayed
majorant.

The sum of its coefficients in all variables other than $B$ is
\[
              1480B^2+\tfrac{887}{2}B+364<2048B^2
                       \qquad(B\ge2).
\]
Every term with testing degree zero contains a factor $\tau$.
For a quadratic form, total absolute coefficient mass bounds the
matrix row norm: an off-diagonal coefficient is split equally
between two matrix entries. Therefore, on the exact support
\[
 I=\{(0,1),(0,2),(1,0),(2,0),(2,1),(4,0),(4,1)\},
\]
\begin{equation}\label{eq:matrixbound}
 Q=Q_0+\sum_I C_{ij}s^iw^j,\qquad
 \norm{Q_0-B_{\rm src}}\le2048B^2\sqrt\eps,
 \qquad \norm{C_{ij}}\le2048B^2.
\end{equation}

The $220$-function organization puts the four already retained
positive terms
\[
 2z[(1+2m)u^2+(1+2D)w^2]N_y^2
       +[zk_1u^2+2zb^2(2D+1)w^2]R^2
\]
into the main matrix. Their additional majorant is
$4B(u^2+w^2)N_y^2+(16Bu^2+4Bw^2)R^2$.
The resulting mass is $1480B^2+943B/2+364<2048B^2$.
Thus \eqref{eq:matrixbound} proves the same guard for this
organization as well. This is a regrouping of four existing squares,
not a new guessed Gram space. The old $z\le1$ row estimate for
$w^2$ is $5$, not $4$:
$2z(1+2D)(1+b)\le4(1+6\eps^2)<5$.

Set $M_R=2048B^2$, $\delta=\delta_R$, and
\[
 T_\delta=2\delta+2\delta^2+\delta^3+\delta^4+\delta^5.
\]
At the stated constants,
\[
 2048B^2\sqrt{\eps_R}\le1/2048,\qquad
 M_RT_\delta\le7M_R\delta=7/512<1/48.
\]
Consequently the explicit source matrices
\[
 H_0^{\rm mat}=Q_0-M_RT_\delta I\succeq I/6,\qquad
 H_{ij}^{\pm}={\delta^{i+j}\over2}(M_RI\pm C_{ij})\succeq0
\]
give the exact split
\[
 Q=H_0^{\rm mat}+\sum_I\left[
 H_{ij}^+\left(1+{s^iw^j\over\delta^{i+j}}\right)
 +H_{ij}^-\left(1-{s^iw^j\over\delta^{i+j}}\right)\right].
\]
Every scalar factor is nonnegative on $0\le s,w\le\delta$.
The circuit remainder has coefficient
\[
 q_C=\tfrac12-(cd_G/S)(2/\lambda+\sqrt x+u).
\]
Its source constant after subtracting $(cd_G/S)\delta^2$ is at least
$1/2-64\eps-4\sqrt\eps-4\delta^2>1/3$. Thus the existing
circuit split with $\sqrt{\delta^2-u}$ is also valid. Source-field
square completion proves the asserted separated SOS.

\section{The nonpositive-angular branch needs no new shear}

Suppose $\lambda\le0$, and put $L=\lambda/\sigma$,
$s_6=\sqrt{1+L^2}$. Nonnegativity of the four terms defining
$\lambda+4pr\sigma$ gives
\[
 H^2\le2pr\le2B,\quad h\le2B\sigma,\quad
 \beta\theta^2\le B\sigma^2,\quad-4B\le L\le0,
 \quad s_6\le5B.
\]
Hence $\theta\ge1/2$, $\beta\theta\le2B\sigma^2$.
(BE3) gives $H\ge H_0/B_s\ge3/4$; therefore
$4\kappa H\ge1$, while $4a,4\ell\ge1$ follows from
\eqref{eq:b4}. Finally (BE4) gives
\[
 d_G={2\mathcal H\over PJ}-{4x\theta\beta\over J\sigma^2}
 \ge {2(1-32B^2\eps)\over1+6B\eps}
                           -{8B\eps\over1-\eps}\ge1.
\]
After multiplying positive denominators, the last inequality is
\[
 1-(64B^2+14B+1)\eps+(16B^2+6B)\eps^2\ge0;
\]
its left side is at least $1-72B^2\eps>0$.

Use the unchanged six-root identity (SR1)--(SR3), on
$v=(N_x,N_y,R,\Gamma,A_\theta,s_6uw(u^2-x)W)^T$.
Its source reference $G_0$ has $N_x$ entry $2$, radial block
$\left(\begin{smallmatrix}2&2b-1\\2b-1&b^2+Q^2\end{smallmatrix}\right)$,
angular entry $4\kappa H$, and $(\Gamma,\widehat V)$ block
$\left(\begin{smallmatrix}d_G&L/(2s_6)\\L/(2s_6)&1\end{smallmatrix}\right)$.
The displayed guards imply $G_0\succeq I_6/3$.

The coefficient matrices printed in (SR2) satisfy, directly,
\[
 \norm{E_{00}}\le128B^2\eps,\qquad
 \norm{E_\gamma}\le128B^2\quad
 (\gamma\in\{10,01,11,20,02,21,30\}).
\]
For an explicit row check, the seven nonconstant rows are bounded,
respectively, by
\[
 49B,\quad64B^2\eps,\quad12B,\quad
 \max(22B,25B^2\eps),\quad8B,\quad8B^2,\quad
 (25/2)B^2\eps.
\]
For $E_{00}$ use the previously proved constant radial estimates
and $s_6x/2\le(5/2)B\eps$.
Thus
\[
 \norm{E_{00}}<1/12,\qquad
 \delta^{|\gamma|}\norm{E_\gamma}\le128B^2\delta
                                      =1/8192<1/56.
\]
The same source matrices (SR6), with $\eta=1/56$ and
$y_\gamma=u^{\gamma_1}w^{\gamma_2}/\delta^{|\gamma|}$, apply
on $0\le u,w\le\delta$, and hence on the common smaller rectangle
$0\le\sqrt u,w\le\delta$. The residual source matrix is
$G_0+E_{00}-I/8\succeq I/8$.
This proves the complementary branch without assuming its guards.
In particular $D>0$ here; its literal zero face occurs in the positive
branch. No inverse of $\beta$ occurs in either construction.

\section{Fixed functions, and what transfer does and does not give}

The positive main family is the twelve pre-Gram functions
\[
 \mathcal V=\{X,Y,Z,W,uZ,wZ,uW,uwW,u^2W,
                 u^{3/2}W,u^{1/2}W,\sqrt{uwWZ_+}\}
\]
times $1$ and the fourteen factors
$\sqrt{1\pm s^iw^j/\delta^{i+j}}$, $(i,j)\in I$.
First apply the fixed angle sign gauge, then the common state sign
making $W\ge0$; here $Z_+=\max(Z,0)$ and $Z_- =\max(-Z,0)$.
Add the seven ordinary functions
\[
 \{wX,uwZ,uX,u^2wW,uY,wY,uwY\},
\]
the nineteen angular functions printed in
\texttt{fixed-features/fixed\_box\_c1.tex}, and its seven circuit
functions $t_i$ together with $\sqrt{\delta^2-u}\,t_i$.
These are $180+7+19+14=220$ labelled functions; the nineteen angular
and seven circuit formulas contain no $R$ or new source choice.
Their verified source-linear root prescriptions remain unchanged.
The six-root family is its ten pre-Gram monomials
\[
 \{X,Y,Z,W,uZ,wZ,uW,uwW,u^2W,u^3wW\}
\]
times $1,\sqrt{1\pm y_\gamma}$, plus the same seven ordinary functions:
$10\cdot15+7=157$.

Alternatively the previously used polynomial testing cover is
\[
 u=\delta^2A_t^8,\quad w=\delta B_t^4,\quad
 \xi^2=1-A_t^4,\quad\eta^2=1-B_t^4,
\]
with the same quarter-state cover and normalization equations.
The existing geometric-sum identities use
exactly the same sixteen multiplier monomials for every fixed $R$.
Its primitive monomial product maps are unchanged. The parent
coefficient embedding contains the newly specified powers of $\delta$;
this statement does not identify it with the previously eliminated
order-$669$ affine source pencil at its old numerical $\delta$.

For rational or algebraic coefficients, one may replace a given real
$R$ by any integer $N\ge R$ and use $\eps_N,\delta_N$. Both constants
are then rational, and every quarter-power used by the cover is algebraic.

All root identities are polynomial after the existing algebraic
covers. Hence testing axes and zero state coordinates are literal
specializations. The two exact angle gauges extend to both signs,
and the two state sectors cover all states. The full (SP6), not a
determinant restriction, is the identity assembled above.

For a bounded positive source $z$ in a real closed extension, some
real $R$ bounds it. Infinitesimal $x,\sigma$ are below this
$\eps_R$, and infinitesimal tests are below this $\delta_R$.
The finite inequalities therefore recover the corresponding
bounded-$z$ infinitesimal admission assertion in B.23 by transfer.
The quantifiers are
\[
 \forall R\ge1\quad\exists(\eps_R,\delta_R,\mathcal V_R)\quad
             \forall\hbox{ native sources and independent tests in the boxes}.
\]
They do not supply one $R$-independent testing space. The displayed
radical branch locations depend on $\delta_R$; a finite cover of the
remaining independent testing regions and source charts is still needed.

\section{An exact obstruction to one proposed fixed-$\delta$ repair}

The need for a new argument is visible in a two-by-two principal
minor. This is an obstruction to the unchanged principal matrix,
not a Bell-slack counterexample or a no-lift theorem.

\begin{proposition}\label{prop:minor}
For every fixed $u_*>0$, arbitrarily large finite $z$ admit native
sources arbitrarily close to the source origin for which the unchanged
twelve-root principal Gram at $(u,w)=(u_*,0)$ is indefinite.
Adding any amount of the positive $J_C^2,J_\ell^2$ radial reserves,
or applying an invertible source-only diagonal congruence, cannot
repair this principal minor.
\end{proposition}

\begin{proof}
It suffices to take $z\ge1$ and
\[
 r=\tfrac12,\quad b={1\over2z},\quad H_0=z+\tfrac12,
 \quad\sigma=t,\quad x=t^2,\quad h=H_0t,
                     \quad0<t\le\eps_z.
\]
Here $B=z+1$. Every strict native gap is positive at this bound:
$h<1$, $E=1-H_0t^2>0$, and $M,J>0$ by \eqref{eq:b1}, or directly
from their formulas. Since $J-t>0$ and $P>1$,
$B_s=P+t^2z(J-t)>1$, so (BE3) gives $\beta>0$ exactly.
The numerator $N_D=K(1+tz)(1+t/2)-JP$ of $D$ satisfies
\[
 N_D=t^2(H_0+\mathcal R),\qquad |\mathcal R|\le19B^4t^2.
\]
Indeed $\mathcal R$ is a polynomial whose monomials all have
$t$-degree at least two, whose $z$-degree is at most four, and whose
absolute coefficient sum is $19$; the checker prints its derivation.
Thus $D>0$. Also
\[
 \lambda\ge2t(H_0^2-p)>0,
\]
because $p\le B$ and $H_0^2-B=z^2-3/4>0$. Finally $c=0$ by
\eqref{eq:b6}, as its reference tends to $2H_0\ge3$ and its error
is bounded explicitly by $1024B^4\eps_z$.

The exact leading coefficients are
\[
 \beta\sim\tfrac12H_0(z+1)t^2,\quad
 D\sim\tfrac12H_0t^2,\quad\lambda\sim4H_0^2t,
\]
\[
 \nu\longrightarrow H_0,\quad d_G,S\longrightarrow2H_0,
                         \quad\rho\longrightarrow0.
\]
At $w=0$, the $(N_x,\Gamma_*)$ principal submatrix is exactly
\[
 \begin{pmatrix}
 2&-2\nu\sqrt{d_G}\,u_*/S\\
 -2\nu\sqrt{d_G}\,u_*/S&1
 \end{pmatrix},
\]
whose determinant tends to $2-2H_0u_*^2$.
For $H_0u_*^2>1$ it is negative for small enough $t$.
An entirely explicit choice for $0<u_*\le1$ is
$z=32/u_*^2$, $0<t\le\eps_z$. The exact formula for $\nu$ gives
$|\nu-H_0|\le16B^2t$; also $\rho=\lambda t\le5Bt$ and
$|S-2H_0|\le128B^2t$. Consequently
$\nu\ge H_0/2$, $d_G\ge H_0$, $S\le3H_0$, and thus
\[
 \det\le2-H_0u_*^2/9<0.
\]
The radial reserves affect only the $(N_y,R)$ rows, and an
invertible congruence preserves inertia. This proves the claim.
\end{proof}

This obstruction must not be overstated. The separately retained
averaged square $(L_0/2)F_A^2$ can change the defective minor:
\[
 F_A=u(Z-a_0wW)={u\over\sigma}\Gamma
                         +{uw\over\sigma kq}T_\theta.
\]
Along the path above at $w=0$, its contribution to the
$\Gamma_*^2$ coefficient tends to $H_0u^2$, precisely the scale
missing from the unchanged principal matrix. A polynomial completion
that reorganizes this square therefore remains open. What fails is
the proposed repair by radial reserves and constant diagonal
normalization alone.

\paragraph{Reproducibility.}
\texttt{check\_bridge.py} uses exact SymPy arithmetic to check the
new positive majorants, their support and coefficient masses, the
two denominator removals and radial divisibility, the reserve
inequalities, the nonpositive $d_G$ guard, and the explicit path's
leading coefficients. Its output is \texttt{check\_results.json}.
It also evaluates the explicit rational source with
$u=2^{-40}$, $z=2^{85}$, and $t=2^{-40}(1+z)^{-8}$.
The native guards and the negative determinant are checked with rational
arithmetic; the source and a rationally congruent $2\times2$ matrix are
printed in \texttt{fixed\_delta\_obstruction.json}.
It does not rerun the $682$ old coefficient equations or regenerate
any old affine pencil. The analytic admissibility proof uses the
native identities throughout, not independent intervals for tied
derived source quantities.
\end{document}
